% !TeX root = ../动手学强化学习笔记.tex
% 由主文档通过 \input 载入；请编译主文档。
\setcounter{chapter}{17} % 保留原书章号。
\chapter{深度模型学习与模型预测控制}
\label{ch:rl-deep-model}

本章将第~\ref{ch:rl-planning} 章的模型思想推广到连续和高维状态，承接函数近似与优化知识，重点讨论动力学预测、模型预测控制和模拟数据的使用。理论分析以简单误差递推为限，学习目标是解释预测误差如何影响规划，并通过真实环境验证决策效果。一个模型不必在所有输入上精确，才能支持有用规划；但平均预测误差很小也不保证决策可靠，因为规划可能专门利用其少数错误区域。

\section{参数化动力学与奖励模型}

确定性预测器 $f_\psi(s,a)$ 输出下一状态，随机模型 $p_\psi(s'\mid s,a)$ 输出分布，奖励模型 $r_\psi(s,a)$ 输出期望奖励。平方误差拟合确定性预测器趋向学习条件均值；多模态后继的均值未必是实际可能状态。负对数似然可训练随机模型，但效果依赖分布族和训练覆盖。

隐空间模型先把观测编码为 $z$，再预测隐状态与奖励。重建观测准确不自动保证保存了决策相关信息；反之，支持控制的表示也不必重建所有像素。评价模型应同时考虑预测、奖励和真实环境中的策略表现。

\section{模型预测控制}

模型预测控制（model predictive control，MPC）在当前状态 $s$ 优化长度 $H$ 的动作序列：
\[
\max_{a_0,\ldots,a_{H-1}}\sum_{k=0}^{H-1}\gamma^k r_\psi(\widehat s_k,a_k)
+\gamma^H\widehat V(\widehat s_H),\qquad
\widehat s_0=s,\quad\widehat s_{k+1}=f_\psi(\widehat s_k,a_k),
\]
并施加合法动作约束。求得序列后只执行第一个动作，观察真实下一状态，再重新优化。随机模型时应优化相应期望或明确的风险目标。随机候选搜索不需要模型可微，梯度规划则要求可微且通常只能找到局部解。

\begin{example}[一步规划的权衡]
系统预测为 $x'=x+a$，动作 $a\in[-1,1]$，当前 $x=1$。若一步目标是最小化 $(x+a)^2+0.1a^2$，求导得 $2+2.2a=0$，所以 $a^*=-10/11$，位于可行区间。它没有直接选 $-1$，因为控制幅度也计入代价。重新规划允许在真实执行后纠正预测偏差。
\end{example}

\section{模型生成数据与误差传播}

学习模型可从真实回放状态出发生成短轨迹，供价值或策略训练使用。起点贴近真实数据有助于控制外推，但模拟样本仍来自同一模型，不应当作独立的新环境证据。真实与模拟数据比例、滚动长度和模型更新频率都影响训练。

设真实确定性动力学为 $f$，两模型在相关区域每步误差至多 $\epsilon$，且 $f$ 对状态为 $L$-Lipschitz。在同一固定动作序列下，预测误差满足 $e_{k+1}\leq Le_k+\epsilon$，故 $e_k\leq\epsilon\sum_{j=0}^{k-1}L^j$。当 $L>1$ 时该界可随时域快速增大。若策略因状态误差而选择不同动作，还需要控制动作敏感性，不能直接沿用这个固定动作界。

\section{不确定性与模型利用}

环境固有随机性即使数据无限也可能存在，模型知识不足则原则上可由新信息减少。模型集成用不同训练样本或初始化的预测差异估计后者，但集成成员可能共同犯错，分歧小不等于真误差小。短时域、真实数据再采样和不确定性惩罚都可减少模型被过度利用的机会，其效果仍需真实交互验证。

\section{习题与实验}
\begin{enumerate}
\item 推导误差递推在 $L<1,L=1,L>1$ 时的界，并解释长期预测难度的差别。
\item 给一维系统模型加入恒定偏差，比较固定开环动作序列与每步重新规划，记录预测误差和真实代价。
\item 构造训练分布下误差很小、在优化器偏爱的动作区域却错误的模型，说明平均验证损失为什么不足以决定规划可靠性。
\end{enumerate}





